\documentclass[10pt,a4paper]{amsart}
\usepackage{amsmath,amssymb,amsthm,graphicx}
\usepackage{verbatim}
\newenvironment{reference}{\comment}{\endcomment}
%\newenvironment{reference}{}{}
\newenvironment{slogan}{\comment}{\endcomment}
\newenvironment{history}{\comment}{\endcomment}
% For commutative diagrams we use Xy-pic
\usepackage[all]{xy}

% We use 2cell for 2-commutative diagrams.
\xyoption{2cell}
% Theorem environments.
%
\theoremstyle{plain}
\newtheorem{theorem}[subsection]{Theorem}
\newtheorem{proposition}[subsection]{Proposition}
\newtheorem{lemma}[subsection]{Lemma}

\theoremstyle{definition}
\newtheorem{definition}[subsection]{Definition}
\newtheorem{example}[subsection]{Example}
\newtheorem{exercise}[subsection]{Exercise}
\newtheorem{situation}[subsection]{Situation}

\theoremstyle{remark}
\newtheorem{remark}[subsection]{Remark}
\newtheorem{remarks}[subsection]{Remarks}

\numberwithin{equation}{subsection}

% Macros
%
\def\lim{\mathop{\mathrm{lim}}\nolimits}
\def\colim{\mathop{\mathrm{colim}}\nolimits}
\def\Spec{\mathop{\mathrm{Spec}}}
\def\Hom{\mathop{\mathrm{Hom}}\nolimits}
\def\Ext{\mathop{\mathrm{Ext}}\nolimits}
\def\SheafHom{\mathop{\mathcal{H}\!\mathit{om}}\nolimits}
\def\SheafExt{\mathop{\mathcal{E}\!\mathit{xt}}\nolimits}
\def\Sch{\mathit{Sch}}
\def\Mor{\mathop{\mathrm{Mor}}\nolimits}
\def\Ob{\mathop{\mathrm{Ob}}\nolimits}
\def\Sh{\mathop{\mathit{Sh}}\nolimits}
\def\NL{\mathop{N\!L}\nolimits}
\def\CH{\mathop{\mathrm{CH}}\nolimits}
\def\proetale{{pro\text{-}\acute{e}tale}}
\def\etale{{\acute{e}tale}}
\def\QCoh{\mathit{QCoh}}
\def\Ker{\mathop{\mathrm{Ker}}}
\def\Im{\mathop{\mathrm{Im}}}
\def\Coker{\mathop{\mathrm{Coker}}}
\def\Coim{\mathop{\mathrm{Coim}}}

% Boxtimes
%
\DeclareMathSymbol{\boxtimes}{\mathbin}{AMSa}{"02}

%
% Macros for moduli stacks/spaces
%
\def\QCohstack{\mathcal{QC}\!\mathit{oh}}
\def\Cohstack{\mathcal{C}\!\mathit{oh}}
\def\Spacesstack{\mathcal{S}\!\mathit{paces}}
\def\Quotfunctor{\mathrm{Quot}}
\def\Hilbfunctor{\mathrm{Hilb}}
\def\Curvesstack{\mathcal{C}\!\mathit{urves}}
\def\Polarizedstack{\mathcal{P}\!\mathit{olarized}}
\def\Complexesstack{\mathcal{C}\!\mathit{omplexes}}
% \Pic is the operator that assigns to X its picard group, usage \Pic(X)
% \Picardstack_{X/B} denotes the Picard stack of X over B
% \Picardfunctor_{X/B} denotes the Picard functor of X over B
\def\Pic{\mathop{\mathrm{Pic}}\nolimits}
\def\Picardstack{\mathcal{P}\!\mathit{ic}}
\def\Picardfunctor{\mathrm{Pic}}
\def\Deformationcategory{\mathcal{D}\!\mathit{ef}}
\usepackage{cleveref}
\crefname{originalreference}{original reference}{original references}
\Crefname{originalreference}{Original reference}{Original references}
\makeatletter
\newlabel{algebra-lemma-fun}{{algebra-lemma-fun}{source}{}{sourcepointer-1}{}}
\newlabel{algebra-lemma-fun@cref}{{[originalreference][0][]algebra-lemma-fun}{[1][0][]source}{}{}{}}
\makeatother
\title{A surjection between finite locally free modules of the same rank is an isomorphism}
\author{The Stacks Project Authors}
\date{stacks-project:modules.tex:local-snapshot-20261012; GNU FDL 1.2 or later}
\begin{document}
\maketitle
\noindent\textit{Editorial note.} All delivered mathematical text is the author's own, copied from the pinned Stacks Project chapter source sources/stacks-modules/pinned-original-source.tex; no statement or proof step was rewritten, reordered or removed.\par\medskip
\noindent\textit{Editorial note.} Preamble provenance: the Stacks Project keeps one single shared preamble for every chapter; this card takes it directly from the staged sources/stacks-modules/preamble.tex as three verbatim spans (lines 8-12, 14-18 and 154-217), so no derived file was written and nothing was added or redefined. The wrapper class and the base packages are supplied by the assembler.\par\medskip
\noindent\textit{Editorial note.} Bibliography provenance: neither the counted statement nor its proof contains a citation (this lane rescans the whole span before writing the card), so the card carries no bibliography and no entry was reconstructed.\par\medskip
\noindent\textit{Editorial note.} Selection provenance: this card was opened under the file-based control contract (CONTROL.md), with the mechanical pre-card gate run on its statement and proof interval before the card was written.\par\medskip
\noindent\textit{Allowed standard prerequisites (editorial).}\begin{itemize}
\item The local triviality of a finite locally free module of a fixed rank. The author's section on locally free sheaves, in which the counted statement occurs. Reduces the statement to a map between finite free modules over the ring of sections of an open subset.
\item The algebraic statement that a surjective endomorphism of a finite free module of fixed rank is an isomorphism. Invoked inside the counted proof at source lines 2386-2388 as a cross-chapter pointer. Upgrades the surjectivity of the map on sections to injectivity, which is what remains to be shown.
\end{itemize}\medskip
\ifdefined\hypertarget\hypertarget{sourcepointer-1}{}\fi
\noindent\textit{Source reference.} Original reference algebra-lemma-fun: sources/stacks-modules/pinned-original-source.tex, source lines 2387--2387. Navigation pointer only: the counted text refers to the Stacks item algebra-lemma-fun, which belongs to another chapter of the project and is not part of this card. The pointer carries the target's own label as its number because this chapter's staged snapshot does not contain the other chapter's tag table; its content is not reproduced here.\par\medskip
\begin{lemma}
\label{lemma-map-finite-locally-free}
Let $(X, \mathcal{O}_X)$ be a ringed space. Let $r \geq 0$.
Let $\varphi : \mathcal{F} \to \mathcal{G}$ be a map of finite
locally free $\mathcal{O}_X$-modules of rank $r$.
Then $\varphi$ is an isomorphism if and only if $\varphi$
is surjective.
\end{lemma}

\begin{proof}
Assume $\varphi$ is surjective. Pick $x \in X$.
There exists an open neighbourhood $U$ of $x$
such that both $\mathcal{F}|_U$ and $\mathcal{G}|_U$ are
isomorphic to $\mathcal{O}_U^{\oplus r}$.
Pick lifts of the free generators of $\mathcal{G}|_U$
to obtain a map $\psi : \mathcal{G}|_U \to \mathcal{F}|_U$
such that $\varphi|_U \circ \psi = \text{id}$.
Hence we conclude that the map
$\Gamma(U, \mathcal{F}) \to \Gamma(U, \mathcal{G})$
induced by $\varphi$ is surjective. Since both
$\Gamma(U, \mathcal{F})$ and $\Gamma(U, \mathcal{G})$
are isomorphic to $\Gamma(U, \mathcal{O}_U)^{\oplus r}$
as an $\Gamma(U, \mathcal{O}_U)$-module we may
apply Algebra, Lemma \ref{algebra-lemma-fun} to see that
$\Gamma(U, \mathcal{F}) \to \Gamma(U, \mathcal{G})$
is injective. This finishes the proof.
\end{proof}

\end{document}
